% Drawing 138/07-14 -- folder 138, batch 07, archive page 131 (PDF page 132).
% Transcription: en marge gauche, au niveau de d) : une parabole coupée par une droite en deux points, un segment à bout marqué ; en dessous, une barre en « T » renversé et une équerre.
% Kind: marginal sketches (parabola and line; T; square). Method 2 (schematic redraw from the facsimile). Status: ok.
% Notes: Left margin at d): a parabola cut by a line in two points, the line continued to a marked end; a stroke crossing the line at the second point; below, an inverted T and a set square (an angle with a crossing stroke). The « -v- o » at the bottom right is text.
% Coordinates: 1 unit = 1 mm = 1/100 of the crop width (bbox in page fractions: [0.112, 0.62, 0.26, 0.705]).
% Redrawn by Claude (claude-opus-5-5), issue #73, 2026-11-07; not checked by a human; not compiled (no TeX in the session).
% Rights: reproduces part of the work -- behind RENDER_TRANSCRIPTS, never published by itself.
\documentclass[tikz,border=3pt]{standalone}
\usepackage{amsmath,amssymb}
\usetikzlibrary{patterns,arrows.meta}
\tikzset{pencil/.style={black!45}, struck/.style={black!35}, mark/.style={circle,draw,inner sep=1.3pt}, pt/.style={circle,fill,inner sep=0.9pt}, opt/.style={circle,draw,fill=white,inner sep=1.1pt}}
\begin{document}
\begin{tikzpicture}[x=1mm,y=1mm,line cap=round,line join=round,font=\scriptsize]
\draw plot[smooth, tension=0.55] coordinates {(18.3,-8.5) (22,-20.7) (32.9,-29.3) (40.2,-22) (43.9,-6.1)};
\draw (14.6,-34.1) -- (75.6,-20.7);
\node[pt] at (32.9,-29.3) {};
\node[pt] at (51.2,-25.2) {};
\node[pt] at (75.6,-20.7) {};
\draw (52.4,-2.4) -- (51.2,-25.2) -- (58.5,-46.3);
\draw[very thick] (6.1,-75.6) -- (41.5,-70.7);
\draw[very thick] (20.7,-61) -- (22,-74.4);
\draw (53.7,-53.7) -- (57.3,-69.5) -- (97.6,-61);
\draw (91.5,-43.9) -- (89,-67.1);
\end{tikzpicture}
\end{document}
